% =====================================================================================
% Abstract. Follows docs/notes/narrative.md (three punchlines). Numbers only from tables/.
% =====================================================================================
System One models, or decision models, answer a question about an input with a calibrated probability instead of
generated text. Within weeks of the first release, in September 2026, they were used as rerankers. We test whether
the decision fine-tuning behind them makes a better reranker than plain reranker fine-tuning of the same model. For
six Qwen backbones from 0.6B to 27B parameters and ten TREC DL and BEIR collections, we compare the untuned model,
asked zero-shot, with its decision fine-tunes and its reranker fine-tunes, two of which we train ourselves. Decision
fine-tuning never makes the better reranker. From 0.6B to 9B, reranker fine-tuning wins on every backbone by 5 to 29
points nDCG@10, also on collections that no fine-tune was trained on, and at 27B the two tie. Up to 9B, decision
fine-tuning adds nothing over simply asking the untuned model. The calibration advantage of decision models depends
on the question. Asked whether a passage answers the query, they are better calibrated than the published reranker
fine-tunes from 4B on. Asked a looser question from prior work, they lose this advantage, and the untuned models gain
as much from the wording. Nor does calibration make up for a worse ranking: up to 9B, a reranker fine-tune keeps the
better passages at a fixed threshold. If you need a reranker, fine-tune a reranker.
